%% LyX 2.4.4 created this file.  For more info, see https://www.lyx.org/.
%% Do not edit unless you really know what you are doing.
\documentclass[english]{article}
\usepackage[T1]{fontenc}
\usepackage[latin9]{inputenc}
\usepackage{xcolor}
\usepackage{booktabs}
\usepackage{calc}
\usepackage{units}
\usepackage{enumitem}
\usepackage{amsmath}
\usepackage{amsthm}
\usepackage{cancel}
\usepackage{geometry}
\geometry{verbose,tmargin=1in,bmargin=1in,lmargin=1in,rmargin=1in}

\makeatletter

%%%%%%%%%%%%%%%%%%%%%%%%%%%%%% LyX specific LaTeX commands.
%% Because html converters don't know tabularnewline
\providecommand{\tabularnewline}{\\}

%%%%%%%%%%%%%%%%%%%%%%%%%%%%%% Textclass specific LaTeX commands.
\numberwithin{equation}{section}
\numberwithin{figure}{section}
\newlength{\lyxlabelwidth}      % auxiliary length 

%%%%%%%%%%%%%%%%%%%%%%%%%%%%%% User specified LaTeX commands.
\usepackage{fancyhdr}
\usepackage{lastpage}
\pagestyle{fancy}
\usepackage{enumitem}
\setlist{noitemsep,nosep}
%\setlist{nolistsep}
\usepackage{ifthen}

%\setlist{nolistsep} % or
%\setlist{noitemsep} %to leave space around whole list

%sets math fonts to be more readable
\everymath=\expandafter{\the\everymath\displaystyle}
%\DeclareMathSizes{10}{10}{10}{10}
\usepackage{multicol}

\usepackage{tikz}
\usetikzlibrary{plotmarks}
\usepackage{pgfplots}
\pgfplotsset{compat=newest}

\usepackage{calc}
\usepackage[nomessages]{fp}

\usepackage{advdate}    % Advancing/saving dates
\usepackage{datetime}   % Dates formatting
\usepackage{datenumber} % Counters for dates
% Added by lyx2lyx
\setlength{\parskip}{\smallskipamount}
\setlength{\parindent}{0pt}

\makeatother

\usepackage{babel}
\begin{document}
\renewcommand{\author}{Dr.\,James Ashe}

\newcommand{\classNum}{132}

\newcommand{\classSec}{C and K}

\newcommand{\nameType}{Solutions}

\newcommand{\examType}{Exam 3}

\renewcommand{\date}{\formatdate{1}{4}{2014}}

%%First page's header%%%%%%%%%%%%%%%%%%%%%%
\fancypagestyle{firststyle} %%header settings for first pagr
{
	\fancyhead[L]{MTH \classNum{}(\classSec{}) \author{}\\
	\textbf{\examType{}} \date{}}
	\fancyhead[C]{\textbf{\nameType}}
	\setlength{\headsep}{14pt}
}
%%%%%%%%%%%%%%%%%%%%%%%%%%%%%%%%%%%%%%%%%%

%%Next page's header%%%%%%%%%%%%%%%%%%%%%%%
\setlist{nolistsep}
\lhead{MTH \classNum{}(\classSec{})} 
\chead{\textbf{\nameType}} 
\cfoot{\thepage{}~of~\pageref{LastPage}} 
\setlength{\headsep}{5pt} 
%%%%%%%%%%%%%%%%%%%%%%%%%%%%%%%%%%%%%%%%%%%

%%Various settings; see the preamble for other settings%%%%%
%\setenumerate[0]{leftmargin=12pt,itemindent=0pt,labelwidth=0pt}
 %\setlist{noitemsep,nosep,topsep=0pt,partopsep=0pt,parsep=0pt,itemsep=0pt}
\renewcommand{\labelenumi}{\textbf{\arabic{enumi}.}}
\renewcommand{\labelenumii}{\textbf{\alph{enumii}.}}
\thispagestyle{firststyle}

%%%%%%%%%%%%%%%%%%%%%%%%%%%%%%%%%%%%%%%%%%
\newcommand{\xmax}{10}
\newcommand{\xmin}{-10}
\newcommand{\xstep}{0} %must be xmin+n
\newcommand{\ymax}{10}
\newcommand{\ymin}{-10}
\newcommand{\ystep}{0} %must be ymin+n

\newcommand{\xspace}{\vspace{1.5cm}}

\emph{You must show work to receive credit.}

\textbf{Find the simple interest AND the future value of the loan. }
\begin{enumerate}[resume]
\item \$1000 at 4.25\% for 6 months.

\begin{enumerate}[resume]
\item [FV:]$1000\left(1+.0425\cdot\frac{6}{12}\right)=\$1021.25$
\item [interest:]$1021.25-1000=\$21.25$\vfill
\end{enumerate}
\end{enumerate}
\textbf{Find the compound amount for the given deposit and the amount
of interest earned.}
\begin{enumerate}[resume]
\item \$2500 at 6.01\% compounded monthly for 38 months.

\begin{enumerate}[resume]
\item [FV:]$2500\left(1+\frac{.0601}{12}\right)^{12\cdot\frac{38}{12}}\approx\$3022.65$
\item [interest:]$3022.65-2500=522.65$\vfill
\end{enumerate}
\item \$2500 at 7\% compounded continuously for $5$ years.

\begin{enumerate}[resume]
\item [FV:]$2500e^{\left(.07\cdot5\right)}\approx\$3547.67$
\item [interest:]$3547.67-2500=\$1047.67$
\end{enumerate}
\end{enumerate}
\textbf{Find the amount that should be invested now to accumulate
the following amount, if the money is compounded as indicated.}
\begin{enumerate}[resume]
\item \$6200 at $9.5\%$ compounded quarterly for 29 months.

\begin{enumerate}[resume]
\item []$6200(1+\frac{.095}{4})^{-(4\cdot\frac{29}{12})}\approx\$4941.41$\vfill
\end{enumerate}
\end{enumerate}
\textbf{Solve}
\begin{enumerate}[resume]
\item Jane needs \$500,000 in 25 years to retire. To meet this goal she
will invest money into an account paying 7\% annual interest rate
compounded continuously. How much money needs to be invested into
the account?

\begin{enumerate}[resume]
\item []$500,000e^{-(25\cdot.07)}\approx\$86,886.97$\vfill
\end{enumerate}
\item You place \$1500 into an account which pays 5\% compounded daily.
How much money will be in the account after 800 days? Assume that
there are 360 days in a year and 30 days in a month.

\begin{enumerate}[resume]
\item []$1500(1+\frac{.05}{360})^{\left(360\cdot\frac{800}{360}\right)}\approx\$1676.27$\vfill\newpage
\end{enumerate}
\end{enumerate}
\textbf{Find the future value of each ordinary annuity, if payments
are made and interest is compounded as given. Then determine how much
of this value is from interest.}
\begin{enumerate}[resume]
\item $R=\$3,250$ (payments); $11\%$ interest compounded semi-annually
for 7 years.\vspace{-5pt}
\end{enumerate}
\begin{itemize}
\item []\begin{multicols}{2}%
\begin{tabular}{rll}
\textbf{N:} & $2\cdot7=14$ & \tabularnewline
${\tt I\%}$: & $11$ & \tabularnewline
${\tt PV}$: & $0$ & \tabularnewline
${\tt PMT}$: & $-3250$ & \tabularnewline
${\tt FV}$: & ? & $\Rightarrow$${\tt FV:}\$65,950.86$\tabularnewline
${\tt P/Y}$: & 2 & \tabularnewline
${\tt C/Y}$: & 2 & \tabularnewline
\end{tabular}\columnbreak
\begin{eqnarray*}
S & = & 3250\left[\frac{\left(1+\nicefrac{.11}{2}\right)^{2\cdot7}-1}{\nicefrac{.11}{2}}\right]\\
 & \approx & \$65,950.86
\end{eqnarray*}
and $65950.86-3250\cdot14=\$20,450.86$ is from interest.\end{multicols}\vfill
\end{itemize}
\textbf{Find the periodic payment that will amount to each given sum
under the given conditions. }
\begin{enumerate}[resume]
\item $S=\$100,000$; interest is 5\% compounded annually; payments are
made at the end of each year for 12 years.\vspace{-5pt}
\end{enumerate}
\begin{itemize}
\item []\begin{multicols}{2}%
\begin{tabular}{rll}
\textbf{N:} & $12\cdot1$ & \tabularnewline
${\tt I\%}$: & 5 & \tabularnewline
${\tt PV}$: & $0$ & \tabularnewline
${\tt PMT}$: & ? & $\Rightarrow$${\tt PMT:-6282.5410}$ \tabularnewline
${\tt FV}$: & $100,000$ & so payments of $\$6282.54$.\tabularnewline
${\tt P/Y}$: & 1 & \tabularnewline
${\tt C/Y}$: & 1 & \tabularnewline
\end{tabular}\columnbreak
\begin{eqnarray*}
R & = & \frac{100000\left(\nicefrac{.05}{1}\right)}{(1+\nicefrac{,05}{1})^{12\cdot1}-1}\\
 & \approx & \$6,282.54
\end{eqnarray*}
\end{multicols}\vspace{-5pt}
\end{itemize}
\textbf{Find the future value of each ordinary annuity, if payments
are made and interest is compounded as given. Then determine how much
of this value is from interest.}
\begin{enumerate}[resume]
\item Payments of $\$800$ into an account paying 6.51\% interest compounded
quarterly for 12 years. \vspace{-5pt}
\end{enumerate}
\begin{itemize}
\item []\begin{multicols}{2}%
\begin{tabular}{rll}
\textbf{N:} & $4\cdot12=48$ & \tabularnewline
${\tt I\%}$: & 6.51 & \tabularnewline
${\tt PV}$: & $0$ & \tabularnewline
${\tt PMT}$: & $-800$ & \tabularnewline
${\tt FV}$: & ? & $\Rightarrow$${\tt FV:}\$57,531.14$.\tabularnewline
${\tt P/Y}$: & 4 & \tabularnewline
${\tt C/Y}$: & 4 & \tabularnewline
\end{tabular}\columnbreak
\begin{eqnarray*}
S & = & 800\left[\frac{\left(1+\nicefrac{.0651}{4}\right)^{4\cdot12}-1}{\nicefrac{..0651}{4}}\right]\\
 & \approx & \$57,531.14
\end{eqnarray*}
Total payments were $800\cdot48=38,400$. So $57,531.14-38,400=\$19,131.14$
is from interest.\end{multicols}
\end{itemize}
\textbf{Find the present value of each ordinary annuity.}
\begin{enumerate}[resume]
\item Payments of \$100 each month for 10 years at 8.5\%.\vspace{-5pt}
\end{enumerate}
\begin{itemize}
\item []\begin{multicols}{2}%
\begin{tabular}{rll}
\textbf{N:} & $12\cdot10=120$ & \tabularnewline
${\tt I\%}$: & 8.5 & \tabularnewline
${\tt PV}$: & ? & \tabularnewline
${\tt PMT}$: & $-100$ & $\Rightarrow{\tt PV:}\$8065.45$\tabularnewline
${\tt FV}$: & $0$  & \tabularnewline
${\tt P/Y}$: & 12 & \tabularnewline
${\tt C/Y}$: & 12 & \tabularnewline
\end{tabular}\columnbreak
\begin{eqnarray*}
P & = & 100\left[\frac{1-\left(1+\nicefrac{.085}{12}\right)^{-10\cdot12}}{\nicefrac{.085}{12}}\right]\\
 & \approx & \$8065.45
\end{eqnarray*}
\end{multicols}\vspace{-5pt}
\end{itemize}
\textbf{Find the monthly house payments necessary to amortize each
loan. Then calculate the total payments and the total amount of interest
paid. }
\begin{enumerate}[resume]
\item \$175,000 at 6.24\% for 30 years.\vspace{-5pt}
\end{enumerate}
\begin{itemize}
\item []\begin{multicols}{2}%
\begin{tabular}{rl}
\textbf{N:} & $12\cdot30=360$\tabularnewline
${\tt I\%}$: & 6.24\tabularnewline
${\tt PV}$: & 175,000\tabularnewline
${\tt PMT}$: & ?\tabularnewline
${\tt FV}$: & $0$ \tabularnewline
${\tt P/Y}$: & 12\tabularnewline
${\tt C/Y}$: & 12\tabularnewline
\end{tabular}\columnbreak
\item [Monthly payments:]${\tt PMT=}-1076.3671$ so payments of $\$1076.37$
\item [Total payments:]$1076.37\cdot360=\$387,493.20$
\item [Interet paid:]$387,493.20-175,000=\$212,493.20$\end{multicols}\vfill
\end{itemize}
\newpage

\textbf{\xspace}

\subsection*{\center{Formula Sheet}}

\noindent\begin{minipage}[t]{1\columnwidth}%
\begin{center}
\begin{tabular}{llll}
$P=$principal/present value & $r=$annual interest rate & $m=$compoundings per year & $n=tm$\tabularnewline
$A=$future/maturity value & $t=$number of years & $i=\nicefrac{r}{m}=$rate per period & $r_{E}=$effective rate\tabularnewline
\end{tabular}
\par\end{center}%
\end{minipage}

\noindent\fcolorbox{black}{white}{\begin{minipage}[t]{1\columnwidth - 2\fboxsep - 2\fboxrule}%
\begin{center}
\begin{tabular}{lclcl}
\textbf{Simple Interest} &  & \textbf{Compound Interest} &  & \textbf{Continuous Compounding}\tabularnewline
\midrule
$A=P(1+rt)$ &  & $A=P(1+i)^{n}$ &  & $A=Pe^{rt}$\tabularnewline
$P=\frac{A}{1+rt}$ &  & $P=A(1+i)^{-n}$ &  & $P=Ae^{-rt}$\tabularnewline
 &  & $r_{E}=\left(1+\frac{r}{m}\right)^{m}-1$ &  & $r_{E}=e^{r}-1$\tabularnewline
\end{tabular}
\par\end{center}%
\end{minipage}}

\xspace%
\noindent\begin{minipage}[t]{1\columnwidth}%
\begin{center}
\begin{tabular}{lll}
$R=$periodic payment & $P=$present value of annuity & $n=$number of periods\tabularnewline
$S=$future value &  & $i=$rate per period\tabularnewline
\end{tabular}
\par\end{center}%
\end{minipage}

\noindent\fcolorbox{black}{white}{\begin{minipage}[t]{1\columnwidth - 2\fboxsep - 2\fboxrule}%
\begin{center}
\begin{tabular}{lcl}
\textbf{Future Value of an Annuity} &  & \textbf{Sinking Fund Payment}\tabularnewline
\midrule
$S=R\left[\frac{\left(1+i\right)^{n}-1}{i}\right]$ &  & $R=\frac{S\cdot i}{(1+i)^{n}-1}$\tabularnewline
\end{tabular}
\par\end{center}%
\end{minipage}}

\noindent\fcolorbox{black}{white}{\begin{minipage}[t]{1\columnwidth - 2\fboxsep - 2\fboxrule}%
\begin{center}
\begin{tabular}{ll}
\textbf{Present Value of an Annuity} & \textbf{Amortization Payments}\tabularnewline
\midrule
$P=R\left[\frac{1-\left(1+i\right)^{-n}}{i}\right]$ & $R=\frac{P}{\left[1-(1+i)^{-n}\right]}$\tabularnewline
\end{tabular}
\par\end{center}%
\end{minipage}}

\xspace

\begin{center}
\textbf{Ti-83/84 }
\par\end{center}

\begin{center}
\begin{tabular}{rl}
\textbf{N} & : number of payments\tabularnewline
${\tt I\%}$ & : interest rate\tabularnewline
${\tt PV}$ & : present value\tabularnewline
${\tt PMT}$ & : payment size\tabularnewline
${\tt FV}$ & : future value\tabularnewline
${\tt P/Y}$ & : payments per year\tabularnewline
${\tt C/Y}$ & : compoundings per year\tabularnewline
\end{tabular}
\par\end{center}

\textbackslash
\end{document}
